% Copyright (c) 2026 CorrectBrain
% SPDX-License-Identifier: MIT

\PassOptionsToPackage{dvipsnames,svgnames,x11names}{xcolor}
\documentclass[tikz,border=8pt]{standalone}
\usetikzlibrary{arrows.meta,positioning,calc,shapes,shadows,decorations.pathmorphing,shapes.geometric,backgrounds,decorations.markings,decorations.pathreplacing,decorations.text,fit,shapes.misc,shapes.arrows,matrix,bending,3d,chains,intersections,shapes.symbols,svg.path,shapes.multipart,snakes,shapes.callouts,angles,quotes}
\providecolor{gold}{RGB}{212,175,55}
\providecolor{silver}{RGB}{192,192,192}
\providecolor{bronze}{RGB}{205,127,50}
\providecolor{darkblue}{RGB}{0,0,139}
\providecolor{lightblue}{RGB}{173,216,230}
\providecolor{darkgreen}{RGB}{0,100,0}
\providecolor{lightgreen}{RGB}{144,238,144}
\providecolor{darkred}{RGB}{139,0,0}
\providecolor{navy}{RGB}{0,0,128}
\providecolor{skyblue}{RGB}{135,206,235}
\providecolor{beige}{RGB}{245,245,220}
\providecolor{cream}{RGB}{255,253,208}
\usepackage{amsmath}
\usepackage{amssymb}
\usepackage{fontawesome5}
\usetikzlibrary{fadings,shadows.blur,patterns,patterns.meta}
\usepackage{xcolor}
% --- CUSTOM THEME COLORS ---
\definecolor{sourcecol}{RGB}{41, 128, 185}   % Belize Hole blue
\definecolor{targetcol}{RGB}{192, 57, 43}    % Pomegranate red
\definecolor{cropcol}{RGB}{127, 140, 141}    % Asbestos gray
\definecolor{textdark}{RGB}{44, 62, 80}      % Midnight blue
\begin{document}
\begin{tikzpicture}[
    textbox/.style={fill=white, draw=gray!40, rounded corners=8pt, inner sep=12pt, blur shadow={shadow opacity=15, shadow yshift=-2pt}}
]
% --- ABSOLUTE BACKGROUND & GRID ---
\fill[white] (-13, -17.5) rectangle (19.2591,15.0183);
\draw[step=1cm, gray!15, very thin] (-13, -17.5) grid (19.2591,15.0183);
% --- ROW 1: TITLE & SUBTITLE ---
\node[font=\bfseries\Huge, textdark, fill=white, inner sep=4pt] at (2.963,13.8664) {The Mechanics of \texttt{object-fit: cover}};
\node[font=\large, gray, fill=white, inner sep=4pt, text width=16cm, align=center] at (3.0243,12.7435) {How the rendering engine scales, crops, and anchors focal points (\texttt{object-position}) in responsive CSS bounding boxes.};
% ==========================================
% STEP 1: THE SOURCE IMAGE (INTRINSIC & FOCAL POINT) [AXIS A: X = -8.0]
% ==========================================
% Row 2 Step Headers (Title baseline: Y = 11.2, Subtitle baseline: Y = 10.8)
\node[font=\Large\bfseries, textdark, anchor=south, fill=white, inner sep=4pt] at (-7.927,10.4208) {1. Source Image};
\node[font=\small, textdark, anchor=north, fill=white, inner sep=2pt] at (-7.927,10.0208) {Intrinsic Size: 600px $\times$ 400px (3:2)};
% Vector Graphic Scene: Sky, Mountains, and Person Silhouette
\begin{scope}
    \clip (-11, 4) rectangle (-5, 8);
    % Sky Gradient
    \shade[top color=skyblue!80!white, bottom color=white] (-11, 4) rectangle (-5, 8);
    % Sun
    \fill[orange!30!yellow!40, opacity=0.5] (-9.8, 6.8) circle (0.9);
    \fill[orange!80!yellow] (-9.8, 6.8) circle (0.45);
    % Mountain Ranges (Depth Layers)
    \fill[sourcecol!35!gray!70] (-11, 4) -- (-10.2, 5.6) -- (-9.0, 4.6) -- (-7.5, 6.0) -- (-6.2, 4.8) -- (-5, 5.8) -- (-5, 4) -- cycle;
    \fill[teal!40!darkgreen!60] (-11, 4) -- (-9.5, 4.9) -- (-8.0, 4.3) -- (-6.0, 5.2) -- (-5, 4.2) -- (-5, 4) -- cycle;
    \fill[green!45!black!70] (-11, 4) to[out=10, in=170] (-5, 4.4) -- (-5, 4) -- cycle;
    % Subject Figure (Head precisely at relative 50% X, 20% Y from Top-Left => (-8, 7.2))
    \begin{scope}[shift={(0, -0.05)}]
        \fill[black!25, blur shadow={shadow opacity=20, shadow yshift=-3pt}]
            (-8.45, 4.0) -- (-8.35, 5.4) to[out=75, in=200] (-8.14, 6.85) -- (-7.86, 6.85) to[out=-20, in=105] (-7.65, 5.4) -- (-7.55, 4.0) -- cycle;
    \end{scope}
    \fill[left color=textdark, right color=textdark!70!black]
        (-8.45, 4.0) -- (-8.35, 5.4) to[out=75, in=200] (-8.14, 6.85) -- (-7.86, 6.85) to[out=-20, in=105] (-7.65, 5.4) -- (-7.55, 4.0) -- cycle;
    \fill[textdark!90!white] (-8.08, 6.85) rectangle (-7.92, 7.02);
    \fill[left color=textdark, right color=textdark!75!black] (-8, 7.2) circle (0.28);
    \fill[textdark!95!black] (-8, 7.28) arc (40:220:0.29) -- cycle;
\end{scope}
% Source Image Border
\draw[sourcecol, ultra thick] (-11, 4) rectangle (-5, 8);
% Focus Crosshair Target at (-8, 7.2)
\draw[red, thick, dashed] (-8, 7.2) circle (0.45);
\draw[red, thick] (-8.6, 7.2) -- (-7.4, 7.2);
\draw[red, thick] (-8, 6.6) -- (-8, 7.8);
\draw[red, thin, dashed] (-8.0, 8.6) -- (-8.0, 7.65);
\node[fill=white, draw=red!70, inner sep=2pt, rounded corners=2pt, font=\sffamily\bfseries\tiny, text=red!85!black] at (-8.0, 8.8) {Focal Point (50\% X, 20\% Y)};
% --- ROW 3: TRANSITION ARROW 1 -> 2 (LOCKED AT Y = 5.2) ---
\draw[-Stealth=3mm, thick, gray] (-5.0, 5.2) -- (-0.5, 5.2) node[midway, above=3pt, font=\bfseries, textdark, fill=white, inner sep=2pt] {Placed into};
% ==========================================
% STEP 2: TARGET CONTAINER (PORTRAIT 1:2) [AXIS B: X = 1.0]
% ==========================================
% Row 2 Step Headers (Title baseline: Y = 11.2, Subtitle baseline: Y = 10.8)
\node[font=\Large\bfseries, textdark, anchor=south, fill=white, inner sep=4pt] at (1.1948,9.8851) {2. Target CSS Box};
\node[font=\small, textdark, anchor=north, fill=white, inner sep=2pt] at (1.1948,9.4851) {Box Size: 300px $\times$ 600px (1:2)};
% Container Box
\fill[targetcol!5] (-0.5, 2) rectangle (2.5, 8);
\draw[targetcol, ultra thick, dashed] (-0.5, 2) rectangle (2.5, 8);
% Dimensional indicators (Fix 1: Gap at Y = 4.6 to 5.7 avoids collision with arrow at Y = 5.2)
\draw[{Stealth[length=4pt]}-Stealth=3mm, gray!80, thin] (-0.5, 1.6) -- (2.5, 1.6) node[midway, below=3pt, font=\scriptsize\sffamily, textdark, fill=white, inner sep=2pt] {300px Width};
\draw[{Stealth[length=4pt]}-, gray!80, thin] (-1.2, 2.0) -- (-1.2, 4.6);
\draw[-{Stealth[length=4pt]}, gray!80, thin] (-1.2, 5.7) -- (-1.2, 8.0);
\node[fill=white, inner sep=2pt, rotate=90, font=\scriptsize\sffamily] at (-1.2, 3.3) {600px Height};
% --- ROW 3: TRANSITION ARROW 2 -> 3 (LOCKED AT Y = 5.2) ---
\draw[-Stealth=3mm, thick, gray] (2.5, 5.2) -- (5.5, 5.2) node[midway, above=3pt, font=\bfseries, textdark, fill=white, inner sep=2pt] {Scale \& Crop};
% ==========================================
% STEP 3: RESULTING CROP (VERTICAL CONTAINER) [AXIS C: X = 10.0]
% ==========================================
% Row 2 Step Headers (Title baseline: Y = 11.2, Subtitle baseline: Y = 10.8)
\node[font=\Large\bfseries, textdark, anchor=south, fill=white, inner sep=4pt] at (10.0487,10.8348) {3. Algorithmic Scale \& Crop};
\node[font=\small, textdark, anchor=north, fill=white, inner sep=2pt] at (10.0487,10.4348) {Resulting Dimensions: 900px $\times$ 600px (1.5x)};
% Scaled Source Artwork (Multiplier = 1.5x)
\begin{scope}
    \clip (5.5, 2) rectangle (14.5, 8);
    \shade[top color=skyblue!80!white, bottom color=white] (5.5, 2) rectangle (14.5, 8);
    % Scaled Sun
    \fill[orange!30!yellow!40, opacity=0.5] (7.3, 6.2) circle (1.35);
    \fill[orange!80!yellow] (7.3, 6.2) circle (0.68);
    % Scaled Mountains
    \fill[sourcecol!35!gray!70] (5.5, 2) -- (6.7, 4.4) -- (8.5, 2.9) -- (10.75, 5.0) -- (12.7, 3.2) -- (14.5, 4.7) -- (14.5, 2) -- cycle;
    \fill[teal!40!darkgreen!60] (5.5, 2) -- (7.75, 3.35) -- (10.0, 2.45) -- (13.0, 3.8) -- (14.5, 2.3) -- (14.5, 2) -- cycle;
    \fill[green!45!black!70] (5.5, 2) to[out=10, in=170] (14.5, 2.6) -- (14.5, 2) -- cycle;
    % Scaled Subject Silhouette (Centered at X = 10.0, Face at Y = 6.8)
    \fill[left color=textdark, right color=textdark!70!black]
        (9.32, 2.0) -- (9.47, 4.1) to[out=75, in=200] (9.79, 6.28) -- (10.21, 6.28) to[out=-20, in=105] (10.53, 4.1) -- (10.68, 2.0) -- cycle;
    \fill[textdark!90!white] (9.88, 6.28) rectangle (10.12, 6.53);
    \fill[left color=textdark, right color=textdark!75!black] (10, 6.8) circle (0.42);
    \fill[textdark!95!black] (10, 6.92) arc (40:220:0.435) -- cycle;
\end{scope}
% The "Crop" Overlay (Semi-transparent white masking overflow)
\fill[white, opacity=0.82] (5.5, 2) rectangle (8.5, 8);   % Left overflow
\fill[white, opacity=0.82] (11.5, 2) rectangle (14.5, 8); % Right overflow
\draw[cropcol, dashed] (5.5, 2) rectangle (8.5, 8);
\draw[cropcol, dashed] (11.5, 2) rectangle (14.5, 8);
% Target Container Outline
\draw[targetcol, ultra thick, dashed] (8.5, 2) rectangle (11.5, 8);
% Scaled Image Outline
\draw[sourcecol, thick] (5.5, 2) rectangle (14.5, 8);
% Centered Focus Crosshair
\draw[red, thick, dashed] (10, 6.8) circle (0.55);
\draw[red, thick] (9.3, 6.8) -- (10.7, 6.8);
\draw[red, thick] (10, 6.1) -- (10, 7.5);
\draw[red, thin] (10.0, 9.2) -- (10.0, 7.35);
\node[fill=white, draw=red!70, inner sep=1.5pt, rounded corners=2pt, font=\sffamily\bfseries\tiny, text=red!85!black] at (10.0, 9.4) {Subject Centered (50\% X)};
% Zero-Overlap Crop Annotations
\node[fill=white, text=cropcol, draw=cropcol, thin, inner sep=2.5pt, rounded corners=2pt, font=\bfseries] at (7.0, 8.5) {Left Crop (300px)};
\node[fill=white, text=cropcol, draw=cropcol, thin, inner sep=2.5pt, rounded corners=2pt, font=\bfseries] at (13.0, 8.5) {Right Crop (300px)};
\node[font=\bfseries, targetcol, fill=white, inner sep=2pt] at (10.0, 1.4) {Visible Area (Center 300px)};
\node[font=\scriptsize\sffamily, textdark!80, fill=white, inner sep=2pt, rounded corners=2pt] at (10.0, 0.6) {\textit{*At 1:2 ratio, Y-overflow is 0px ($\text{Scale}_H$ dictates). The 20\% Y-position is dormant here.}};
% ==========================================
% STEP 4: RESPONSIVE HORIZONTAL STRETCH (THE 50% 20% FOCUS LOCK) [AXIS D: X = -2.0]
% ==========================================
\node[font=\bfseries\Large, textdark, fill=white, inner sep=4pt] at (-2.0, -1.5) {4. The 50\% 20\% Focus Point Lock (Responsive Horizontal Container)};
\node[font=\small, textdark, fill=white, inner sep=2pt] at (-2.0, -2.2) {Container: 800px $\times$ 300px (8:3) $\implies$ Width Dictates Scale ($1.333\times$), Activating the Vertical 20\% Y-Crop};
\begin{scope}[yshift=-2.0cm]
    % Container: (-6.0, -6.0) to (2.0, -3.0) -> Width = 800px (8.0), Height = 300px (3.0)
    % Scaled Image: Width = 800px (8.0), Height = 533.3px (5.333)
    % Y-Overflow = 233.3px (2.333). Top Crop = 46.7px (0.467), Bottom Crop = 186.7px (1.867)
    % Scaled bounds: X from -6.0 to 2.0; Y from -7.867 to -2.533
    \begin{scope}
        \clip (-6.0, -7.867) rectangle (2.0, -2.533);
        \shade[top color=skyblue!80!white, bottom color=white] (-6.0, -7.867) rectangle (2.0, -2.533);
        % Scaled Sun
        \fill[orange!30!yellow!40, opacity=0.5] (-4.4, -4.13) circle (1.2);
        \fill[orange!80!yellow] (-4.4, -4.13) circle (0.6);
        % Scaled Mountains
        \fill[sourcecol!35!gray!70] (-6.0, -7.867) -- (-4.93, -5.73) -- (-3.33, -7.07) -- (-1.33, -5.2) -- (0.4, -6.8) -- (2.0, -5.47) -- (2.0, -7.867) -- cycle;
        \fill[teal!40!darkgreen!60] (-6.0, -7.867) -- (-4.0, -6.67) -- (-2.0, -7.47) -- (0.67, -6.27) -- (2.0, -7.6) -- (2.0, -7.867) -- cycle;
        \fill[green!45!black!70] (-6.0, -7.867) to[out=10, in=170] (2.0, -7.33) -- (2.0, -7.867) -- cycle;
        % Scaled Subject Silhouette (Head at X = -2.0, Y = -3.60 => exactly 20% below container top!)
        \fill[left color=textdark, right color=textdark!70!black]
            (-2.60, -7.867) -- (-2.47, -6.0) to[out=75, in=200] (-2.19, -4.07) -- (-1.81, -4.07) to[out=-20, in=105] (-1.53, -6.0) -- (-1.40, -7.867) -- cycle;
        \fill[textdark!90!white] (-2.11, -4.07) rectangle (-1.89, -3.85);
        \fill[left color=textdark, right color=textdark!75!black] (-2.0, -3.60) circle (0.373);
        \fill[textdark!95!black] (-2.0, -3.50) arc (40:220:0.387) -- cycle;
    \end{scope}
    % Top and Bottom Overflows Masked
    \fill[white, opacity=0.82] (-6.0, -2.533) rectangle (2.0, -3.0);
    \fill[white, opacity=0.82] (-6.0, -6.0) rectangle (2.0, -7.867);
    \draw[cropcol, dashed] (-6.0, -2.533) rectangle (2.0, -3.0);
    \draw[cropcol, dashed] (-6.0, -6.0) rectangle (2.0, -7.867);
    % Target Container and Scaled Image Outlines
    \draw[targetcol, ultra thick, dashed] (-6.0, -6.0) rectangle (2.0, -3.0);
    \draw[sourcecol, thick] (-6.0, -7.867) rectangle (2.0, -2.533);
    % Crosshair at Locked Focal Point (-2.0, -3.60)
    \draw[red, thick, dashed] (-2.0, -3.60) circle (0.48);
    \draw[red, thick] (-2.6, -3.60) -- (-1.4, -3.60);
    \draw[red, thick] (-2.0, -4.2) -- (-2.0, -3.0);
    \draw[red, thin, dashed] (-2.0, -1.8) -- (-2.0, -3.1);
    \node[fill=white, draw=red!70, inner sep=2pt, rounded corners=2pt, font=\sffamily\bfseries\tiny, text=red!85!black] at (-2.0, -1.6) {Focal Point Locked at 20\% Y of Container};
    % Dimensional Callouts for Step 4 (Fix 2: Relocated to the left side at X = -6.3)
    \draw[{Stealth[length=4pt]}-Stealth=3mm, thick, cropcol] (-6.3, -2.533) -- (-6.3, -3.0);
    \node[anchor=east, text width=3.5cm, align=right, font=\scriptsize\bfseries, cropcol, fill=white, inner sep=2pt] at (-6.45, -2.766) {Top Crop: 46.7px ($20\% \times \Delta Y$)};
    \draw[{Stealth[length=5pt]}-Stealth=3mm, thick, targetcol] (-6.3, -3.0) -- (-6.3, -6.0);
    \node[anchor=east, text width=3.5cm, align=right, font=\scriptsize\bfseries, targetcol, fill=white, inner sep=2pt] at (-6.45, -4.5) {Visible Height: 300px};
    \draw[{Stealth[length=4pt]}-Stealth=3mm, thick, cropcol] (-6.3, -6.0) -- (-6.3, -7.867);
    \node[anchor=east, text width=3.5cm, align=right, font=\scriptsize\bfseries, cropcol, fill=white, inner sep=2pt] at (-6.45, -6.933) {Bottom Crop: 186.7px ($80\% \times \Delta Y$)};
    \draw[{Stealth[length=4pt]}-Stealth=3mm, thick, targetcol] (-6.0, -8.3) -- (2.0, -8.3) node[midway, below=4pt, font=\scriptsize\bfseries, targetcol, fill=white, inner sep=2pt] {Target Box Width: 800px};
    % Step 4 Mathematical Callout Box (Fix 3: Relocated to Axis C at X = 10.0)
    \node[textbox, text width=9.5cm] at (10.0, -5.0) {
        \textbf{\large \faIcon{crosshairs} \ Focal Point Alignment Math (\texttt{50\% 20\%})}\\[4pt]
        When horizontal aspect ratio causes vertical overflow ($\Delta Y > 0$), the Y-offset formula actively shifts the image up to anchor the subject's face:\\[3pt]
        $\bullet$ \textbf{Scale Factor}: $\max\left(\frac{800}{600}, \frac{300}{400}\right) = \max(1.333, 0.75) = \mathbf{1.333x}$\\[2pt]
        $\bullet$ \textbf{Scaled Size}: $800\text{px} \times 533.3\text{px} \implies \Delta Y = 533.3 - 300 = \mathbf{233.3px}$\\[2pt]
        $\bullet$ \textbf{Top Crop} = $\Delta Y \times 20\% = 233.3\text{px} \times 0.20 = \mathbf{46.7px}$\\[2pt]
        $\bullet$ \textbf{Bottom Crop} = $\Delta Y \times 80\% = 233.3\text{px} \times 0.80 = \mathbf{186.7px}$\\[3pt]
        \textit{Result: Face remains at exactly 20\% from the container top!}
    };
\end{scope}
% ==========================================
% ROW 4: EXPLANATORY TEXT BOXES (BOTTOM FOOTER - STRICT TOP-ALIGNED AT Y = -11.5)
% ==========================================
% Box 1: Scale Calculation & Directives [Axis D: X = -2.0]
\node[textbox, anchor=north, text width=11.5cm] at (-2.0, -11.5) {
    \textbf{\Large \faIcon{expand-arrows-alt} \ The Core Directive of \texttt{cover}}\\[4pt]
    The algorithm must scale the image so its dimensions are $\ge$ the container's dimensions, strictly preserving the aspect ratio. The dimension requiring the \textit{larger} multiplier dictates the global scale factor:\\[4pt]
    $\text{Scale}_W = \frac{300}{600} = 0.5 \qquad \text{Scale}_H = \frac{600}{400} = 1.5$\\[4pt]
    $\text{Final Scale Factor} = \max(0.5, 1.5) = \textbf{1.5x}$
};
% Box 2: Default Crop Calculation [Axis C: X = 10.0]
\node[textbox, anchor=north, text width=9.0cm] at (10.0, -11.5) {
    \textbf{\Large \faCrop \ The Default 50\% Crop}\\[4pt]
    By default, \texttt{object-position} is \texttt{50\% 50\%}. This mathematically centers the image, dividing the overflow equally:\\[4pt]
    $\bullet$ Scaled Width = $600\text{px} \times 1.5 = \textbf{900px}$\\
    $\bullet$ Total Overflow = $900\text{px} - 300\text{px} = \textbf{600px}$\\
    $\bullet$ Left / Right Crop = $600\text{px} \div 2 = \textbf{300px}$
};
\end{tikzpicture}
\end{document}
